% 第 17 章　一个振子叫振动，一群振子叫波
\chapter{一个振子叫振动，一群振子叫波}

\step{反常识开场}

\begin{mainline}
体育课上，你站在操场这头，朝一百米外的小明喊他的名字。大约三分之一秒后，他回过头来了。

现在请你认真想一个问题：你的声音是怎么过去的？

最省事的回答是“空气把声音送过去的”。可是如果真的有一股空气从你的嘴边出发，以每秒三百四十米的速度飞向小明，那应该是一阵狂风——足够吹乱他的头发、掀翻他的作业本。事实上什么也没有被吹过去。小明身边那团空气一直待在小明身边，你嘴边这团空气也一直待在你嘴边。

没有东西过去，却有东西到了。

类似的怪事还有很多。在游泳池的一端拍一下水面，另一端的塑料鸭子会上下晃动几下——但它几乎不朝你漂过来，也不漂走。体育场里几万名观众玩“人浪”，浪头几秒钟绕场一周，可每个人都好好坐在自己的座位上，谁也没有离开过。浪明明“到了”，人却一步也没有动。

这一章就是要搞清楚：那个“没有东西过去却确实到了”的东西，到底是什么。
\end{mainline}

\step{先猜一猜}

\begin{mainline}
在往下读之前，先押个注。找一根长绳子的画面在脑子里放一遍：绳子一头拴在门把手上，你捏着另一头，手腕快速上下抖一下，一个“鼓包”沿着绳子跑向门把手。

请你对下面三句话各给一个直觉判断：

\begin{itemize}
  \item 鼓包跑得有多快？它和你的手抖得多用力有关，还是和绳子绷得多紧有关？
  \item 鼓包到达门把手以后，会怎样？消失、停住，还是掉头跑回来？
  \item 如果你和另一个人各捏绳子一头，同时各抖一下，两个鼓包在绳子中间迎面相遇，它们会像两个皮球一样撞开、弹回去吗？
\end{itemize}

把你的答案记下来。这一章结束时，我们会回来批改。我可以先剧透一句：绝大多数人在第三题上会答错，而答错的原因恰恰是最有用的那一课——它会逼你把“波”和“会飞的物体”在脑子里彻底分家。
\end{mainline}

\step{动手实验}

\begin{experiment}[一根绳子上的波]
你需要：一根三到五米长的绳子（跳绳、晾衣绳、背包带都行），一个同伴，一部手机。

第一步，定性观察。把绳子一头拴在桌腿上，你捏住另一头，让绳子刚好离开地面、略微绷直。手腕快速向上甩一下再收回。你会看到一个鼓包贴着绳子跑出去，撞上桌腿，然后——仔细盯着——它会掉头跑回来。这就是你在上一节押过注的第二题。

第二步，让鼓包“立住”。把绳子绷得更紧一些，手腕以稳定的节奏左右（或上下）连续摆动，像打着拍子。调整节奏，你会找到某个特定的快慢，让绳子不再乱晃，而是呈现出一个稳定的、中间最鼓的形状，像被定格的波浪。再加快一点节奏，还能找到让绳子分成两段、三段反向摆动的节奏。记住“节奏对了形状才稳”这个感觉，本章后半段它有个正式的名字：驻波。

第三步，定量测一次。让同伴在绳子那头猛拉一下制造一个鼓包，你用手机秒表测鼓包从一头跑到另一头的时间，测三次取平均。用卷尺量出绳长，波速就是绳长除以时间。然后把绳子绷得更紧，重测一次。你会发现：同一根绳子，绷得越紧，鼓包跑得越快。波速不取决于你抖得多用力，而取决于绳子本身的状态。这一条线索我们会在“数学翻译器”里兑现成一个公式。

安全提示：别用弹力很强的橡皮绳，绷紧后脱手会抽到人。
\end{experiment}

\step{旧模型遇到麻烦}

\begin{mainline}
到现在为止，这本书教你描述世界的方式是：找到一个个物体，写下每个物体的位置、速度和受力，然后用牛顿定律推算它们下一步怎样运动。弹簧振子、单摆都是这么处理的。那么，拿这套办法来对付绳子上的鼓包，会发生什么？

麻烦一：物体太多了。鼓包经过的地方，绳子并没有一小段“搬家”到桌腿那边去——实验里看得清清楚楚，鼓包跑过之后，每一小段绳子都还留在原来的位置附近，只是上下动了一下。所以“鼓包”根本不是一块移动的物体。那它是谁？你总不能对着“什么东西都没有”画受力图。水波把这层尴尬放得更明显：池塘里的落叶被浪经过时只是原地起伏，可浪头明明一路推进了十几米。如果坚持“运动必须有运动者”，我们就得承认一个荒诞结论——有一个看不见的东西跑过了水面。

麻烦二：就算我们退而求其次，把绳子切成一百小段，给每一小段都当作一个振子来追踪，立刻碰到一个尴尬：每一小段为什么非得跟着邻居动不可？第一段绳子自己并没有被你的手碰到，它凭什么动起来？旧模型里，一个物体运动状态的改变需要受合力；这里每一段绳子受的力，来自它两边邻居的拉扯。也就是说，想让牛顿定律继续管用，我们必须先承认一件事——这一百个振子不是一百个互不相干的振子，它们彼此之间拴在一起，会互相传递“该动了”的消息。

麻烦三：能量去哪了的账对不上。你的手抖一下，只做了一瞬间的功，随后手就停了。可鼓包带着能量一路跑到桌腿，撞上的时候甚至能让桌腿轻轻响一声。能量确确实实搬了家，搬家的载体却不是任何一块物质。

三条麻烦指向同一个出路：我们需要一个新概念，它描述的不是“哪个物体在哪里”，而是“哪个状态传到了哪里”。
\end{mainline}

\step{发明新概念}

\begin{mainline}
像历史上的物理学家一样，我们给这个新概念起个名字：\textbf{波}。

先把绳子的新形象固定下来：一根绳子，可以看成一串紧紧挨着的振子，每一个振子都被邻居用弹性力拴着。你抖动手边这一端，第一个振子偏离平衡位置；它一偏，就拉扯第二个振子；第二个被拉起来之后再去拉扯第三个……每个振子只做一件事：在自己家附近原地振动，顺便推一把下一家。所谓“鼓包跑过去”，跑的不是绳子，而是“偏离平衡”这个\textbf{状态}，以及跟着状态一起走的能量。

这里值得请出最常用的比喻：体育场的人浪。它像的地方在于——每个人都只在原地站起坐下，浪头却能绕场一周；传播的是“起立”这个状态，不是观众。它不像的地方在于——人浪靠每个人用眼睛看、用脑子反应，所以波速由人的反应时间决定，而且可以被随意指挥；绳子上的波靠的是物理的弹性力，波速完全由绳子的松紧和粗细决定，你一旦松手，谁也指挥不动它。记住这个“不像”，它能拦住很多想当然。

按照振子振动的方向，波分成两大类：

\begin{itemize}
  \item \textbf{横波}：振子原地振动的方向，与波前进的方向垂直。绳子上的鼓包就是横波——绳子上下动，波水平走。
  \item \textbf{纵波}：振子振动的方向，与波前进的方向平行。把弹簧放在地上，在一头快速一推一拉，你会看到一截“被挤紧”的部分沿着弹簧跑过去。空气里的声音就是纵波：空气分子只是在原地前后挤让，把“密一点、疏一点”的状态一站一站传下去。这就是开场喊话的正确图像——传到小明耳朵里的不是你的空气，而是空气“挤了一下”这个消息。
\end{itemize}

还有一个对以后至关重要的观察：波速由谁决定？实验已经给了线索——同一根绳子绷得越紧，波越快。换成更粗更重的绳子呢？直觉（以及后面的公式）会告诉你：更重的绳子每一小段更“懒”，被邻居拉起来更费劲，波就更慢。也就是说，\textbf{波速由介质自己的性质（绷得多紧、材料多重多硬）决定，而与波源的振幅无关}。你的手抖得再用力，也只是把鼓包造得更大，并不能让它跑更快。

最后强调一遍波到底运走了什么。振子本身原地不动，却有两样东西实实在在地抵达了远方：一样是\textbf{能量}——你的手做的功，变成鼓包撞响桌腿的那一声；放大到自然尺度，海啸在开阔大洋上每小时推进八百公里，浪高不过一米，几乎不被察觉，可它携带的能量抵达浅海时能把浪抬到十几米高，把码头的船拍上岸。另一样是\textbf{信息}——鼓包的形状、节奏、先后次序，都是可以被对岸读取的“消息”。守岛的渔民看浪的形状判断远处是否有风暴，靠的就是波忠实地把扰动的细节运了过来。能量与信息能脱离物质单独旅行，这是波最深刻、也最容易被一句“水过去了”掩盖掉的性质。
\end{mainline}

\begin{figure}[htbp]
\centering
\begin{tikzpicture}[scale=1]
  \draw[->] (-0.2,0) -- (8.6,0) node[right] {波前进方向};
  \draw[->] (0,-1.5) -- (0,1.8) node[above] {振子偏移};
  \draw[thick,domain=0:8,samples=120] plot (\x,{sin(\x*90)});
  \foreach \p in {0.5,1.5,2.5,3.5,4.5,5.5,6.5,7.5}
    \fill (\p,{sin(\p*90)}) circle (1.6pt);
  \draw[<->] (1,1.25) -- (5,1.25) node[midway,above] {一个波长 $\lambda$};
  \draw[->,thick] (2.5,{sin(2.5*90)}) -- (2.5,{sin(2.5*90)+0.75});
  \node at (3.6,-1.1) {黑点只上下振动，并不随波右移};
\end{tikzpicture}
\caption{横波：波形向右传播，每个振子（黑点）只在原地上下振动。}
\end{figure}

\begin{figure}[htbp]
\centering
\begin{tikzpicture}[scale=1]
  \foreach \i in {0,...,48}
    \fill ({0.17*\i + 0.22*sin(\i*30)},0) circle (1.3pt);
  \node at (1.35,0.6) {疏部};
  \node at (3.1,0.6) {密部};
  \node at (4.85,0.6) {疏部};
  \node at (6.6,0.6) {密部};
  \draw[->,thick] (7.9,0.5) -- (8.9,0.5) node[right] {波前进方向};
  \node at (4.2,-0.7) {纵波：振子沿波的方向前后挤让，疏密状态在传播};
\end{tikzpicture}
\caption{纵波：每个点仍在原处附近，传播的是“疏”与“密”的状态。}
\end{figure}

\step{一句话模型}

\begin{oneline}
波是一群互相拴着的振子在玩接力：每个振子只在原地振动，把“偏离平衡”的状态和能量一站一站传下去；传走的是消息，不是物质。
\end{oneline}

\begin{mainline}
这句话里每个词都值钱。“互相拴着”回答了为什么第二个振子会动；“只在原地”回答了为什么没有物质远行；“一站一站”回答了为什么波有确定的速度——接力的快慢由每一棒的交接速度决定，也就是由介质决定；“消息”这个词先欠着，下一章讲声音与信息时我们会连本带利收回来。背下这句话，本章后面的每一个公式都只是给它加注脚。
\end{mainline}

\step{数学翻译器}

\begin{mainline}
现在把直觉翻译成数字和图像。先解决最实用的三个量。

\textbf{周期与频率}来自第 16 章：每个振子上下一次的时间叫周期 \(T\)，每秒振动次数叫频率 \(f=1/T\)，单位是赫兹。波源抖多快，整条链上每个振子就抖多快——频率是波源的印记。

\textbf{波长}是空间里的问法：在任何一个瞬间拍张快照，相邻两个鼓包峰顶之间的距离，记作 \(\lambda\)。

\textbf{波速} \(v\) 是传播的快慢。这三个量被一个朴素的关系拴在一起。想一想：波源每完成一次振动（用去时间 \(T\)），就多造出一个完整的鼓包；而队伍最前面的鼓包已经跑出去一段距离 \(vT\)。所以新造出来的峰和前面的峰之间恰好相隔 \(vT\)，这就是波长：

\[
  \lambda = vT, \qquad \text{也就是} \qquad v = f\lambda .
\]

这个式子描述什么？它描述同一种波里频率与波长的此消彼长。为什么这样写？因为“振一次、造一个峰、峰间距被波速锁定”这三件事逼着它长这样。何时成立？只要波速 \(v\) 在介质里处处一样——同一根均匀绳子、同一团静止空气里都成立。何时失效？当介质本身在不同位置或对不同频率给出不同波速时（比如水有深有浅、光穿过玻璃），\(v\) 不再是常数，就要对每种频率分开处理。

用极端情形检查一下：频率加倍，波长减半——波源抖得越快，峰越挤，合理；波速趋于零的介质里，任何频率都传不出去，波长挤成零——介质“冻住”了，也合理。
\end{mainline}

\begin{mainline}
\textbf{两种图像别搞混。}描述波时有两张完全不同的图，初学者最常在这里摔跤。第一张是\textbf{波形图}：某一瞬间给整根绳子拍张快照，横轴是位置 \(x\)，纵轴是该处振子的偏移——它回答“此刻哪里鼓、哪里瘪”。第二张是\textbf{振动图}：盯住某一个振子不放，给它一个人拍录像，横轴是时间 \(t\)，纵轴是它的偏移——它回答“这一个点什么时候上、什么时候下”。前者的峰间距是波长（空间的尺子），后者的峰间距是周期（时间的尺子），两者用波速换算。一个常用的读图本领：已知波形图和波的传播方向，判断某个振子下一瞬间往哪动——把整条波形沿传播方向挪一小截，看那个位置上的新高度是更高还是更低即可。波形整体平移，而每个点只在自己的竖线上滑动，这句话就是“传走的是状态，不是物质”的图像版。
\end{mainline}

\begin{mainline}
\textbf{再给耳朵算一笔账。}空气中的声速约 \(340\) 米每秒，人耳能听到的频率范围大约从 \(20\) 赫兹到 \(20000\) 赫兹。由 \(\lambda = v/f\) 换算成波长：最低的音对应约 \(17\) 米——比一间教室还长；最高的音对应约 \(1.7\) 厘米——比一粒纽扣还小。这个跨度直接解释了一批日常经验：低音鼓的波长和墙壁、门窗的尺寸相当，很容易绕过去，所以楼下装修的“咚咚”声无孔不入；而钥匙碰撞的高音波长只有几厘米，遇到障碍基本沿直线传播，躲在柱子后面就明显变闷。声音三要素里的音调、响度、音色分别对应频率、振幅和波形形状，那是下一章的主题，这里只需记住：它们全都是可以用 \(v=f\lambda\) 换算的实在尺寸。
\end{mainline}

\begin{mainline}
\textbf{一次真实的估算。}实验里你测过绳上波速，现在算它该是多少。对一根张紧的绳，分析每一小段绳元受的侧向拉力（推导放进数学桥），可以得到

\[
  v = \sqrt{\frac{T}{\mu}},
\]

其中 \(T\) 是绳子的张力，\(\mu\) 是单位长度绳子的质量（线密度）。取一根真实的吉他弦做例子：中张力钢弦张力约 \(70\) 牛，一米长弦的质量约 \(5\) 克，即 \(\mu = 0.005\) 千克每米。代入：

\[
  v = \sqrt{\frac{70}{0.005}} = \sqrt{14000} \approx 118\ \text{米每秒}.
\]

这个数字马上有用。吉他弦长约 \(0.65\) 米，后面讲驻波时你会知道，一根两头固定的弦，最基本的振动是半个波长刚好卡在弦长里，即 \(\lambda = 2\times 0.65 = 1.3\) 米。于是它最基本的频率是

\[
  f = \frac{v}{\lambda} = \frac{118}{1.3} \approx 91\ \text{赫兹},
\]

离真实吉他第五弦的空弦音（110 赫兹）只差一个“我们选了根偏松的弦”的距离。一根弦、一个公式、一个数量级吻合的音高——这就是“波速由介质决定”这句话可以拿出手的证据。
\end{mainline}

\begin{mathbridge}
上面说“分析每一小段绳元”具体怎么做？设绳上偏离平衡的距离为 \(y\)，它同时是位置 \(x\) 和时间 \(t\) 的函数 \(y(x,t)\)。取一小段绳子，两边张力方向的微小差别提供指向平衡位置的回复合力；这段绳子的质量是 \(\mu\,\mathrm{d}x\)，加速度是 \(y\) 对时间的二阶导数。把牛顿第二定律写上去、取小振幅近似（波形坡度处处很小），两边消去 \(\mathrm{d}x\)，得到

\[
  \frac{\partial^2 y}{\partial t^2} = \frac{T}{\mu}\,\frac{\partial^2 y}{\partial x^2}.
\]

这就是一维波动方程。读它的方式是：左边是“这一处振子的时间加速度”，右边是“这一处波形的空间弯曲程度”乘以一个由介质决定的常数。整句话翻译过来：\textbf{某处波形弯得越厉害，那里的振子被拉回平衡位置的加速度就越大}。时间变化率与空间弯曲度被拴在一起，这正是“振动会传播”的数学写法。可以验证 \(y(x,t)=A\sin(kx-\omega t)\) 是它的解，条件是 \(\omega/k=\sqrt{T/\mu}=v\)——即 \(v=f\lambda\) 穿上了正弦函数的外衣。其中 \(k=2\pi/\lambda\) 叫波数，\(\omega=2\pi f\) 叫角频率。小振幅近似什么时候失效？当波形坡度不再小，比如巨浪翻卷、弦被猛拨到出现折角，方程右边就得补上更复杂的项，下面“模型的边界”还会回来收这笔账。
\end{mathbridge}

\begin{mainline}
\textbf{两列波相遇：本章最重要的反直觉事实。}回到“先猜一猜”的第三题。两个鼓包迎面相遇时，真实发生的是：在相遇的那一小段绳子上，两处的偏移直接相加——相遇瞬间绳子该在哪，就是两列波各自说“该在哪”的和。然后，两列波像什么也没发生一样，各自穿过对方，带着原来的形状和速度继续走。它们不像皮球，不会像台球那样撞开。

这叫\textbf{叠加原理}：同一介质中同时存在几列波时，任一振子的实际偏移等于各列波单独引起的偏移之和。它是数学桥里那个方程的直接推论——方程里只出现 \(y\) 本身而没有 \(y\) 的平方、三次方，所以“解加解还是解”。

叠加原理最亲近的证据在你的耳朵里：一间屋子里两个人同时说话，两列声波在空气中穿过彼此，你却能同时听清两个人在说什么——如果声波像车辆一样会“相撞”，声音早就乱成一锅粥了。也正因如此，耳朵和麦克风收到的永远是叠加后的总振动，能从总振动里把两个人分开，靠的全是大脑（或算法）的事后功夫。波在介质里各自独立传播，这是它和粒子最不一样的地方之一。

叠加原理一口气结出三颗果子：

\begin{itemize}
  \item \textbf{干涉}：两个同频率的波源（比如两根手指以同一节奏点水面）在水面上造出两圈波纹，叠加后有些地点两峰总相遇、振动格外强，有些地点一峰一谷总相遇、水面几乎不动。强弱的分布只由“从两个波源到那一点的路程差”决定。
  \item \textbf{反射与折射}：鼓包跑到桌腿那头发觉“过不去了”，能量没地方去，只好掉头——这就是实验里看到的反射，而且固定端会把鼓包翻个面（上鼓包回来变下鼓包）。如果波从一种介质走进波速不同的另一种介质——比如水波从深水区走进浅水区，波速变慢而频率不变，由 \(v=f\lambda\) 波长被迫变短，波的前进方向也随之偏折——这就是折射。海边波浪总是差不多平行地涌上沙滩，正是因为近岸水浅波慢，浪头被“掰”向岸边。
  \item \textbf{衍射}：波遇到障碍物或小孔时，不会死板地沿直线投出几何阴影，而是会绕到障碍物的侧后方去。缺口越小（与波长相当时），绕得越明显。你在门外能听见屋里人说话，很大一部分靠的就是声波绕过门框；而说话声里低音（波长长）比高音（波长短）更容易“拐过弯”，所以隔墙听到的声音总是闷闷的。
\end{itemize}
\end{mainline}

\begin{mainline}
\textbf{驻波：被边界关起来的波。}一根两头固定的弦上，向右跑的波在端点反射成向左跑的波，两列频率相同、方向相反的波叠加，绝大多数频率会叠加成一团乱晃的噪声；只有少数几个特殊频率，能让弦上恰好整数个“半波长”严丝合缝地卡进弦长，叠加出一个形状稳定、原地起伏的波形：有些点永远不动（波节），有些点振动最猛（波腹）。这就是驻波，你在实验第二步亲手找到过它。允许的频率是

\[
  f_n = n\,\frac{v}{2L}, \qquad n=1,2,3,\dots
\]

其中 \(L\) 是弦长。这一串离散的频率叫这根本弦的\textbf{本征频率}——弦的“身份证号码”。弦乐器（吉他、小提琴、钢琴）靠它定音高：琴码和琴枕固定 \(L\)，调音拧的旋钮改的是张力 \(T\)，也就是改 \(v\)。管乐器是同一出戏换了舞台：笛子、小号里的空气柱是纵波振子链，管长替代弦长，吹奏者改变指法就是在改变空气柱的有效长度，换一组本征频率。你在厨房就能验证：往玻璃瓶里吹瓶口，会听到一个低沉的音；往瓶里加些水再吹，音调明显升高——水面上升缩短了空气柱，\(L\) 变小，由 \(f_n = nv/(2L)\) 频率自然抬高。为什么你随便敲桌子是“咚”一声，而琴弦能唱出稳定的音？因为桌子形状复杂、允许的驻波频率挤在一起不成比例，琴弦的本征频率却整齐地排成 \(1:2:3:\cdots\)，耳朵把这种整齐听成了“音高”。
\end{mainline}

\begin{figure}[htbp]
\centering
\begin{tikzpicture}[scale=0.95]
  \begin{scope}
    \draw (0,0) -- (6,0);
    \draw[thick,domain=0:6,samples=60] plot (\x,{0.7*sin(\x*30)});
    \node[right] at (6.2,0) {$n=1$，半个波长};
  \end{scope}
  \begin{scope}[yshift=-2cm]
    \draw (0,0) -- (6,0);
    \draw[thick,domain=0:6,samples=60] plot (\x,{0.7*sin(\x*60)});
    \node[right] at (6.2,0) {$n=2$，一个波长};
  \end{scope}
  \begin{scope}[yshift=-4cm]
    \draw (0,0) -- (6,0);
    \draw[thick,domain=0:6,samples=90] plot (\x,{0.7*sin(\x*90)});
    \node[right] at (6.2,0) {$n=3$，一个半波长};
  \end{scope}
\end{tikzpicture}
\caption{两头固定的弦上允许存在的三种驻波。只有整数个半波长能恰好卡进弦长。}
\end{figure}

\step{模型的边界}

\begin{mainline}
一个诚实的模型要说清自己在哪里下班。

\textbf{边界一：需要介质——本章的波全是机械波。}绳子、水面、空气柱、大地，波都是“一群物质振子的接力”。没有物质的地方，接力棒传不出去。这不是模型的缺陷，而是它的定义；但它立刻逼出一个新问题：阳光穿过真空来到地球，那是什么在接力？这个悬念先挂账，本书电磁篇会回答：电场与磁场本身就是空间的物理属性，它们的变化按麦克斯韦方程的要求互相牵制着向前传播，不需要任何物质介质。到那时你会庆幸本章先学会了“状态传播”这个抽象动作——电磁波只是把它从物质身上搬到了场身上。

\textbf{边界二：小振幅与线性。}叠加原理、\(v=f\lambda\)、波动方程，全都建立在“偏移很小、回复力与偏移成正比”的假设上。浪足够高时会翻卷破碎，声音足够强时会压成激波（音爆），那时波速开始依赖振幅，强波会追上并吞掉前面的弱波，叠加原理失效，波形一边传一边变形。日常说话、琴弦、池塘涟漪都在安全区里；爆炸、音爆、碎浪在区外。

\textbf{边界三：介质均匀、静止，波源也静止。}如果介质自己在动（风、水流），波的表观速度是介质运动与波速的叠加；如果介质不均匀（海水温度分层、大气上冷下热），波速逐点变化，波会被连续地“掰弯”。声波在夏日午后贴着热地面传播时会向上弯，声音“传不远”，夜深人静时反而能听到远处的火车——不只是夜里更安静这么简单，而是冷空气层把声波掰回了地面。还有一种越界是波源自己在动：救护车的笛声逼近时变尖、远去时变低，这不是声速变了，而是波源一边跑一边吐波纹，把前方的波长挤短、后方的波长拉长。这个现象叫多普勒效应，它是下一章的正式成员，这里只登记一下它的户口本在本章：波长确实被挤变了，但 \(v=f\lambda\) 里的 \(v\) 始终由空气说了算。

\textbf{典型误用}也要点名。一，“波把东西运过去了”——错，运走的是状态与能量，浮子和观众都留在原地。二，“频率由介质决定”——错，频率跟着波源走，介质只决定波速，波长是被二者挤出来的结果。三，“驻波是两列波停下来了”——错，是两列波仍在以原速对穿，只是叠加出的轮廓恰好站着不动。
\end{mainline}

\begin{challenge}
同一个介质里，不同频率的波跑得一样快吗？对理想的绳子和空气，答案是一样快——波动方程里的 \(v\) 不含频率。但对深水波、玻璃里的光，答案是各频率速度不同，这叫\textbf{色散}。色散一旦成立，就得区分两种速度。\textbf{相速度} \(v_p=\omega/k\)，是波形上某一固定相位点（比如某个峰）的前进速度；\textbf{群速度} \(v_g=\mathrm{d}\omega/\mathrm{d}k\)，是一整包波（波包）整体的移动速度，能量和信息都跟着波包走，所以物理上“信号的速度”指群速度。在一个水塘里可以亲眼看到两者的差别：扔一块石头进去，盯着波圈外缘看，你会发现单个波纹从波圈的内侧生成，穿过波圈，在外缘消失——单波比整个波圈跑得快。谁快谁慢，由介质的色散关系 \(\omega(k)\) 决定。到了量子篇你会再见到它们：物质波的相速度可以超过光速而群速度不能，正是群速度对应粒子运动的速度，所以因果律安然无恙。
\end{challenge}

\step{回头解释开场现象}

\begin{mainline}
现在批改开场作业。

喊话为什么没有风？因为声波是纵波接力：你嘴边的空气被声带推着前后挤让，每个空气分子只在原地前后移动不到一根头发丝的距离，把“密—疏—密”的状态一站一站传给邻居，\(0.3\) 秒后状态抵达小明的耳朵，推动他的耳膜。空气一个没走，消息一条没落。频率由你的声带决定（这就是小明能认出“是你的声音”的原因），波速 \(340\) 米每秒由空气的温度和弹性决定，波长被二者挤出。

泳池里的鸭子为什么只上下晃？因为水波里的每个水分子（在理想情形下）只在原地绕小圈，波传的是起伏的状态和能量，鸭子是水面的浮标，跟着原地画圈，自然不漂走——除非有风或水流在真正搬动水体，那是另一回事。

人浪为什么能绕场？它是“状态接力”最赤裸的演示：站起坐下的是人，传播的是“轮到你了”。还可以顺手做个估算：体育场上人浪扫过一整圈大约二十几秒，跑道一圈约四百米，浪头速度大约每秒十几米；相邻观众间隔约半米，也就是说“起立”状态每秒交接二十几次，和一个人看到邻座起身再做出反应所需的零点几秒完全对得上。人浪的速度表里没有张力和质量，只有反应时间——这正是“波速由介质决定”在一个奇特“介质”上的翻版。

还有一件开场没提、但值得用新模型收割的事：地震预警。地震同时发出两种体波——纵波（P 波，约 \(6\) 千米每秒）和横波（S 波，约 \(3.5\) 千米每秒），横波破坏力大得多。监测台网先收到跑得快、破坏小的 P 波，抢在 S 波到达之前用电磁波（速度三十万千米每秒，约等于瞬间）向更远的城市发警报。成都的地震预警系统在多次地震中为市区争取到几十秒：高铁减速、手术暂停、电梯停靠最近楼层。这几十秒全部来自本章的一个核心事实：\textbf{不同介质里，波速由介质决定，而且彼此差得很远}。
\end{mainline}

\step{脑洞实验与挑战题}

\begin{thought}[地月钢绳]
想象在地球和月球之间拉一根绷直的钢绳（先不计较它会因为自重断成多少截）。你在地球这头用大锤敲一下，月球那头的人多久能知道？

用本章模型算：钢是固体，纵波在钢里的速度约 \(5000\) 米每秒，地月距离约 \(38\) 万千米，敲击的“挤压状态”要接力传过去，耗时约 \(380000000/5000 = 76000\) 秒，也就是大约 \(21\) 个小时。而无线电波走同样的路只要 \(1.3\) 秒。同样是“敲一下、传个信”，机械波和电磁波差了近六万倍。

再想一层：月球上的宇航员把耳朵贴在绳子上能听见，飘在绳子旁边的宇航员却什么都听不见——真空中没有空气振子链，声波无处接力。科幻电影里太空中惊天动地的爆炸声，属于导演，不属于物理。
\end{thought}

\begin{mainline}
最后，把本章连回全书的地图。我们从一个振子（第 16 章）走到一群振子（本章），靠的是“邻居之间有耦合”这一条。这个配方威力极大，后面会一再复用：把“一群振子”换成“一排电荷”，耦合换成电场磁场，就走向电磁波；把琴弦的本征频率离散谱记住，量子力学里电子在原子中同样只允许驻波式的一组离散状态——能量量子化的最早直觉就来自“整数个半波长卡进有限空间”；甚至现代场论把空间每个点都看作一个振子，粒子是场的振动量子。一个振子叫振动，一群振子叫波——这句话的下半句，要等你读到量子篇才听得到。
\end{mainline}

\begin{puzzles}
  \item 宇航员在空间站舱内敲击舱壁，舱内的同伴能听到声音；舱外太空行走的同伴（没贴着头盔）却听不到。分别说明两次“听到”与“听不到”的接力链各是什么。（提示：舱内一条链是金属舱壁加舱内空气，舱外缺了哪一环？）
  \item 吉他手拧紧琴弦，音调升高。请用本章公式从“拧旋钮”推到“音变高”，写出推理链的每一步，不需要数字。（提示：旋钮改的是哪个量？它先影响 \(v\) 还是 \(f\)？弦长变了吗？）
  \item 山谷里对着远处崖壁喊一声，一会儿听到回声。有人说“声波撞到崖壁被撞回来了”。这个说法对了一半错了一半：反射确实是能量掉头，但“撞”字掩盖了边界条件的区别。绳波在固定端反射会翻面（峰变谷），你猜声波在坚硬的崖壁上反射会不会翻面？（提示：想一想崖壁处的空气分子能不能自由移动，它更像绳子的固定端还是自由端。）
  \item 体育场人浪的“波速”由什么决定？和绳上波速的决定因素相比，哪个更像“介质的弹性”，哪个更像“介质的惯性”？如果你让观众们反应再快一倍，人浪会变快吗——这对应绳子上改变什么？（提示：人浪没有物理回复力，它的“接力速度”由观察—反应链决定；想想这个类比在“波速由谁决定”这一点上能走多远。）
\end{puzzles}
